Consider a 3D Constitutive law relating the second Piola-Kirchhoff stress tensor \(\bm{S}\) to Green's Lagrangian strain tensor \(\SolidStrain\).

Recall the Frobenius inner product for two second order tensors \(\bm{A}\) and \(\bm{B}\) with coordinate representations \(A^{ij} \triangleq \mathbf{G}^i\cdot \bm{A}\mathbf{G}^j\) and  \(B_{ij} \triangleq \mathbf{G}_i\cdot \bm{B}\mathbf{G}_j\) on a basis \(\{\mathbf{G}_k\}\)
\begin{equation*}
    \left<\bm{A},\,\bm{B}\right> = \operatorname{tr} \bm{A}\bm{B}^{\mathrm{t}} = A^{ij} B_{ij}
\end{equation*}
If \(\delta \SolidStrain = \left(\hat{\bm{B}} \delta\FrameStrain\right) \stackrel{s}{\otimes}\FrameBasis\) for some second-order tensor \(\hat{\bm{B}}\), then
\begin{equation*}
    \left<\bm{S},\,\delta\SolidStrain\right> = \SectionIntegral{ 
    \operatorname{tr} \left[\bm{S}\frac{1}{2}\left(\hat{\bm{B}}\delta \FrameStrain \otimes \FrameBasis + \FrameBasis\otimes \hat{\bm{B}}\delta \FrameStrain \right)
    \right]
    }
\end{equation*}
Exploiting linearity of the trace together with the identity \(\operatorname{tr} (\bm{a}\otimes\bm{b}) = \bm{a}\cdot \bm{b}\) for arbitrary vectors \(\bm{a}\) and \(\bm{b}\) it follows
\begin{equation*}
    \left<\bm{S},\,\delta\SolidStrain\right> = \delta \FrameStrain\cdot \SectionIntegral{
        \hat{\bm{B}}^{\mathrm{t}}\frac{1}{2}\left(\bm{S}\FrameBasis + \bm{S}^{\mathrm{t}}\FrameBasis\right)
    }
\end{equation*}
and consequently by requiring \(\left<\FrameResult,\,\delta\FrameStrain\right>=\left<\bm{S},\,\delta\SolidStrain\right>\)
\begin{equation*}
\FrameResult \equiv \SectionIntegral{\hat{\bm{B}}^{\mathrm{t}}\frac{1}{2}\left(\bm{S}\FrameBasis + \bm{S}^{\mathrm{t}}\FrameBasis\right)}
\end{equation*}
For symmetric \(\bm{S}\) one further has
\begin{equation}
\FrameResult = \SectionIntegral{
  \hat{\bm{B}}^{\mathrm{t}} \bm{S}\FrameBasis
}
\end{equation}
This has the intuitive implication that ony the first ``column'' and/or ``row'' of the stress tensor \(\bm{S}\) should contribute to the frame stress resultants \(\ShearResult\) and \(\CurveResult\).

The stress tensor can be expressed as follows
\begin{equation}
\bm{S}=\sigma \FrameBasis \otimes \FrameBasis+\bm{\tau} \otimes \FrameBasis+\FrameBasis \otimes \bm{\tau}
\label{eq:saint-venant-stress}
\end{equation}
\paragraph{Linearization}
\begin{equation*}
    D {\FrameResult} [\delta \FrameStrain] = \hat{\bm{B}}^{\mathrm{t}}D \bm{S} [D \SolidStrain [\delta \FrameStrain]]
\end{equation*}
Using \(\mathbb{C}[\delta \SolidStrain] \triangleq D\bm{S}[\delta \SolidStrain]\) for vectors \(\delta \SolidStrain\) tangent at \(\SolidStrain\) 
\begin{equation*}
 D {\FrameResult} [\delta \FrameStrain] = \hat{\bm{B}}^{\mathrm{t}} \mathbb{C} D \SolidStrain [\delta \FrameStrain]
\end{equation*}
Now using \(D \SolidStrain [\delta \FrameStrain] = (\hat{\bm{B}}\delta \FrameStrain) \stackrel{s}{\otimes} \FrameBasis\)
\begin{equation}
\hat{\mathbf{K}} = \SectionIntegral{
    \hat{\mathbf{B}}^{\mathrm{t}} \mathbb{C}\hat{\bm{B}} + \hat{\bm{G}}
    }
\label{eq:section-tangent}
\end{equation}

\subsection{Truncation}
When the truncated strain \cref{eq:wagner-gl-strain} is used, it is convenient to rewrite \cref{eq:section-tangent} as the sum of explicit \textbf{Linear} \(\hat{\bm{K}}_{\mathrm{L}}\) and \textbf{Wagner} \(\hat{\bm{K}}_{\mathrm{W}}\) terms given by
\begin{equation}
\hat{\bm{K}}_{\mathrm{L}} = \SectionIntegral{
\left[\begin{matrix}
\mathbb{C} & - \mathbb{C} \bm{r}^{\times} & \mathbb{C} \FrameBasis\otimes \bm{\varphi} & \mathbb{C} \FrameBasis_{\beta} \otimes \nabla_{\beta} \bm{\varphi}\\
%\bm{r}^{\times} \mathbb{C} 
& - \bm{r}^{\times} \mathbb{C} \bm{r}^{\times} 
& \bm{r}^{\times} \mathbb{C} \FrameBasis\otimes \bm{\varphi} & \bm{r}^{\times} \mathbb{C} \FrameBasis_{\beta} \otimes \nabla_{\beta} \bm{\varphi}\\
% \bm{\varphi} \otimes \FrameBasis \mathbb{C} 
& % - \bm{\varphi} \otimes \FrameBasis \mathbb{C} \bm{r}^{\times} 
& \bm{\varphi} \otimes \FrameBasis \mathbb{C} \FrameBasis\otimes \bm{\varphi} 
& \bm{\varphi} \otimes \FrameBasis \mathbb{C} \FrameBasis_{\beta} \otimes \nabla_{\beta} \bm{\varphi}\\
\mathrm{sym.} % \nabla_{\beta} \bm{\varphi} \otimes \FrameBasis_{\beta} \mathbb{C} 
& % - \nabla_{\beta} \bm{\varphi} \otimes \FrameBasis_{\beta} \mathbb{C} \bm{r}^{\times} 
& % \nabla_{\beta} \bm{\varphi} \otimes \FrameBasis_{\beta} \mathbb{C} \FrameBasis\otimes \bm{\varphi} 
& \nabla_{\beta} \bm{\varphi} \otimes \FrameBasis_{\beta} \mathbb{C} \FrameBasis_{\beta} \otimes \nabla_{\beta} \bm{\varphi}
\end{matrix}\right]
}
\label{eq:section-linear-tangent}
\end{equation}
and using \(r^2 \triangleq \bm{r}\cdot\bm{r}\)
\begin{equation}
\begin{aligned}
\hat{\bm{K}}_{\mathrm{W}} =
\kappa_T & \SectionIntegral{
\left[\begin{matrix}
\bm{0} & r^{2}  \mathbb{C} \FrameBasis \otimes \FrameBasis & \bm{0} & \bm{0} \\
r^{2}  \FrameBasis \otimes \FrameBasis \mathbb{C} 
& r^{2}  \left(\bm{r}^{\times} \mathbb{C} \FrameBasis \otimes \FrameBasis -  \FrameBasis \otimes \FrameBasis \mathbb{C} \bm{r}^{\times}\right) & r^{2}  \FrameBasis \otimes \FrameBasis \mathbb{C} \FrameBasis\otimes \bm{\varphi} & r^{2}  \FrameBasis \otimes \FrameBasis \mathbb{C} \FrameBasis_{\beta} \otimes \nabla_{\beta} \bm{\varphi}\\
\bm{0} & r^{2}  \bm{\varphi} \otimes \FrameBasis \mathbb{C} \FrameBasis \otimes \FrameBasis & \bm{0} & \bm{0} \\
\bm{0} & r^{2}  \nabla_{\beta} \bm{\varphi} \otimes \FrameBasis_{\beta} \mathbb{C} \FrameBasis \otimes \FrameBasis & \bm{0} & \bm{0} 
\end{matrix}\right]} \\
+\kappa_T^2 & \SectionIntegral{\left[\begin{matrix}
\bm{0} & \bm{0} & \bm{0} & \bm{0} \\
\bm{0} & r^{4} \FrameBasis \otimes \FrameBasis \mathbb{C} \FrameBasis \otimes \FrameBasis & \bm{0} & \bm{0} \\
\bm{0} & \bm{0} & \bm{0} & \bm{0} \\
\bm{0} & \bm{0} & \bm{0} & \bm{0}
\end{matrix}\right]
} \\
&+\SectionIntegral{\left[\begin{matrix}
\bm{0} & \bm{0} & \bm{0} & \bm{0} \\
\bm{0} & r^2 \sigma & \bm{0} & \bm{0} \\
\bm{0} & \bm{0} & \bm{0} & \bm{0} \\
\bm{0} & \bm{0} & \bm{0} & \bm{0} 
\end{matrix}\right]
}
\end{aligned}
\label{eq:section-wagner-tangent}
\end{equation}

\subsection{Change of Reference}
\begin{equation*}
    \bm{u}_A(\bm{r}_A) = \bm{u}_A(\bm{0}) + \bm{\theta}\times\bm{r}_A + \bm{\alpha}\cdot\bm{\varphi}_A(\bm{r}_A) \,\FrameBasis
\end{equation*}
%
\begin{equation*}
\bm{u}_A(\bm{0}) = \bm{u}_B(\bm{0}) + \bm{\theta}\times \bm{\SectionLocation_{AB}}
\end{equation*}
%
Thus
\begin{equation*}
\begin{pmatrix}
\delta\bm{\gamma}_A \\ \delta \bm{\kappa}_A
\end{pmatrix}
= \begin{bmatrix}
    \bm{1} & -\bm{\SectionLocation_{AB}^{\times}} \\
    \bm{0} & \bm{1}
\end{bmatrix}
\begin{pmatrix}
\delta \bm{\gamma}_B \\ \delta \bm{\kappa}_B
\end{pmatrix}
\end{equation*}


\subsection{Summary}